\ulohahead{specializace UI-SU \textnormal{\footnotesize[úroveň 3]}}{POMDP a aktualizace beliefu}
\begin{zadani}
Částečně pozorovatelný MDP se dvěma stavy $s_1,s_2$. Pro akci $a$: $T(s_1\!\to\!s_1)=0{,}8$, $T(s_1\!\to\!s_2)=0{,}2$, $T(s_2\!\to\!s_1)=0{,}3$, $T(s_2\!\to\!s_2)=0{,}7$. Model pozorování: $O(o\mid s_1)=0{,}9$, $O(o\mid s_2)=0{,}2$. Počáteční belief $b=(0{,}5,\,0{,}5)$.
\begin{enumerate}[label=(\alph*)]
  \item \bod{1} Definujte POMDP a pojem belief (víra o stavu). \emph{(Toto je rozsah požadavku: ,,POMDP - základní definice''.)}
  \item Napište vzorec aktualizace beliefu po akci $a$ a pozorování $o$. \emph{(volitelně, nad rámec požadavku)}
  \item Spočtěte nový belief po provedení $a$ a pozorování $o$. \emph{(volitelně, nad rámec požadavku)}
\end{enumerate}
\end{zadani}

\begin{reseni}
\textbf{(a)} \emph{POMDP} $=(S,A,T,R,\Omega,O,\gamma)$: jako MDP, ale agent stav přímo \emph{nevidí}; po akci dostane jen \emph{pozorování} $o\in\Omega$ s pravděpodobností $O(o\mid s',a)$ (obecně $O:\Omega\times S\times A\to[0,1]$ - pozorování smí záviset i na akci). \emph{Belief} $b$ je rozdělení pravděpodobnosti nad stavy (co agent ví); rozhoduje se podle něj (POMDP lze chápat jako MDP nad spojitým prostorem beliefu).

\textbf{(b)} Aktualizace (předpověď přes přechod + korekce pozorováním, Bayes):
\[
b'(s')\ \propto\ O(o\mid s',a)\sum_{s} T(s'\mid s,a)\,b(s),
\]
poté normalizace, aby $\sum_{s'}b'(s')=1$.

\textbf{(c)} (V tomto příkladu $O$ na akci nezávisí, tj.\ $O(o\mid s')$.) \emph{Predikce}: $\tilde b(s_1)=0{,}8\cdot0{,}5+0{,}3\cdot0{,}5=0{,}55$, $\tilde b(s_2)=0{,}2\cdot0{,}5+0{,}7\cdot0{,}5=0{,}45$.
\emph{Korekce} pozorováním $o$: $\propto(0{,}9\cdot0{,}55,\ 0{,}2\cdot0{,}45)=(0{,}495,\ 0{,}09)$. Normalizace ($/0{,}585$): $b'(s_1)\approx0{,}846$, $b'(s_2)\approx0{,}154$.
\end{reseni}

\begin{jadro}
POMDP = MDP + skrytý stav + pozorování $O(o\mid s',a)$; agent drží belief $b$ (rozdělení nad stavy). Update: $b'(s')\propto O(o\mid s')\sum_s T(s'\mid s,a)b(s)$, pak normuj. (Stejná logika jako filtrace HMM s akcemi.)
\end{jadro}

\kalibrace{Požadavek chce u POMDP \emph{jen základní definici} (část a) - to je povinné jádro. Belief update (b, c) je nad rámec požadavku; uvádíme ho pro pochopení (je to táž logika jako filtrace HMM), ale na zkoušce se nevyžaduje.}
\past{Belief update = predikce (přechod) PAK korekce (pozorování) + normalizace, jako u HMM. Belief je rozdělení, ne jeden ,,nejpravděpodobnější'' stav.}
